\subsection{What the Distributional Stress-Test Reveals}

Every model in this study achieved AUROC 1.00 on the internal UCI test set. By the standard metrics that dominate published clinical ML literature, all five would be described as performing excellently. When the same models were applied to the MIMIC-IV demo cohort, AUROC fell to values indistinguishable from chance (0.48--0.58), ECE exceeded 0.68 for every model, and conformal coverage dropped from near-target to 0.21--0.25.

Before interpreting these numbers, the stress-test framing matters. The MIMIC demo is an open-access 100-patient subset released for pipeline development and workshop use, not a formally designed validation cohort. Its CKD prevalence (23.7\%) differs from the UCI training population (62.5\%) by 39 percentage points, and seven of 24 model features are entirely absent from the demo data. The purpose of including it is to illustrate what happens when models encounter a population that does not resemble their training data, not to estimate clinical deployment performance.

That said, the pattern is exactly what calibration theory predicts. Isotonic recalibration that achieved ECE 0.000 internally did not survive the shift. Echouffo-Tcheugui and Kengne found external calibration validation to be the exception rather than the norm in CKD prediction literature~\cite{EchouffoTcheugui2012}. Liou and colleagues documented a similar gap in a deployed malnutrition prediction model, finding calibration drift and subgroup bias that required post-deployment recalibration within a large healthcare system~\cite{Liou2024}. This study provides a concrete, reproducible illustration of why that gap matters: perfect internal metrics say nothing about behaviour under distribution shift.

\subsection{What Drove the External Failure}

Three factors contributed to the calibration collapse. The first is prevalence shift. A model trained and calibrated on a 62.5\% CKD population assigns probabilities near the training prevalence. The MIMIC cohort sits at 23.7\% CKD, so those probability estimates are systematically too high, producing calibration curves that fall above the diagonal at every predicted probability level. This is the most direct driver of high MIMIC ECE.

The second factor is feature missingness. Seven of the 24 features in the UCI schema are not routinely recorded in MIMIC. Urine-specific gravity, urine sugar, pus cells, pus cell clumps, bacteria, appetite, and pedal edema were all imputed using UCI training-set medians and modes. Those imputed values carry no information about individual MIMIC patients. From the model's perspective, seven features are effectively noise columns on the external cohort.

The third factor is dataset saturation. AUROC 1.00 on a 61-patient test set, with near-perfect cross-validation scores, signals that the UCI benchmark does not offer a realistic discrimination challenge. The learned decision boundaries are sharp enough to separate every patient in the training domain. Those boundaries do not generalize to a population with a different disease severity distribution and different measurement patterns.

XGB was the least poor performer externally (AUROC 0.579, ECE 0.680). That small margin appears to reflect gradient boosting's regularization limiting overfitting to UCI-specific patterns. Even so, the margin is not large enough to change the practical conclusion.

\subsection{Conformal Prediction as a Diagnostic Tool}

The conformal coverage collapse offers something calibration numbers alone do not: a direct, interpretable signal of distribution shift. Conformal prediction has been applied in clinical settings precisely because it provides this kind of interpretable coverage guarantee, including recent applications to individual-level diagnostic uncertainty in chronic disease~\cite{Vazquez2022,Sreenivasan2025}. The theoretical guarantee of split conformal prediction holds only when calibration and test data come from the same distribution~\cite{AngelopoulosBates2022}. Coverage of 0.22--0.25 on MIMIC, against a 0.90 target, is a quantitative statement that the MIMIC population is not exchangeable with the UCI validation set. A system operator who sees conformal coverage fall from 0.97 to 0.22 gets an immediate signal that the model is operating outside its valid domain.

That is the practical value of including conformal prediction in a deployment framework, separate from its theoretical properties. It creates a coverage audit trail. High MIMIC singleton rates for LR (1.00) and NB (0.99) might appear to suggest interpretable outputs. They do not: when coverage is only 0.24, a singleton prediction is a confident wrong answer for roughly three quarters of patients. Singleton rate without coverage context is not a useful interpretability metric.

\subsection{Comparison to Published CKD Models}

The eight externally-validated CKD models identified by Echouffo-Tcheugui and Kengne reported a range of calibration outcomes, but assessment was typically informal, often through visual calibration plots rather than ECE or Brier Score computation~\cite{EchouffoTcheugui2012}. More recent models, including the KFRE validated by Tangri and colleagues~\cite{Tangri2016} and its UK validation by Major and colleagues~\cite{Major2019}, achieve strong external discrimination (C-statistic 0.80--0.90) in prospective cohort studies with purpose-built feature ascertainment.

The AUROC values seen on MIMIC in this study (0.48--0.58) are not comparable to those results, because the KFRE and similar models were validated on cohorts where inputs were actually measured. The comparison instead reinforces a narrower point: a model applied to harmonized data with seven imputed features is not the same as a model applied to complete data. Deployment performance depends on what the deployment environment actually records, not on what the training environment provided.

\subsection{Limitations}

The most important limitation to state plainly: the MIMIC cohort used here is the publicly available 100-patient demonstration subset released by PhysioNet for pipeline development, not a formally designed external validation cohort. It was not selected to represent a specific clinical population, it contains only 23 CKD cases, and seven of the features were entirely missing and had to be imputed from UCI training statistics. The stress-test framing in this paper is intentional and accurate. The MIMIC results show what distribution shift looks like numerically; they do not estimate performance in any real clinical deployment.

At 97 patients after filtering, calibration metrics carry wide uncertainty regardless of the above. Bootstrap confidence intervals for MIMIC ECE span as much as 0.18 for some models. Sample size requirements for precise external calibration assessment substantially exceed the 97 patients available in this stress-test cohort~\cite{Riley2024EvalPart3}. The point estimates should be read as directional, not precise.

Feature harmonization made the missingness problem worse. The seven imputed features were filled using UCI training-set statistics, so those columns carry no information specific to individual MIMIC patients. That is an inherent constraint of cross-dataset harmonization when features are domain-specific, and it cannot be corrected analytically after the fact.

Subgroup analysis was additionally limited by sample size. Diabetes and hypertension subgroups each fell below the 10-patient exclusion threshold, and the age subgroups (below 65: $n$\,=\,55; 65 and above: $n$\,=\,42) were small enough that bin-level ECE estimates carry meaningful uncertainty. Any subgroup findings should be treated as exploratory.

\subsection{Future Directions}

Three areas appear worth pursuing. First, the full MIMIC-IV database would provide several thousand patients with CKD-relevant laboratory data, allowing larger-scale calibration assessment and properly powered subgroup analysis. The extraction pipeline from this study is directly applicable to the full database.

Second, the deployment readiness framework here is intentionally simple. Criteria are binary and thresholds are set a priori without formal sample size calculations. Extending the framework to account for uncertainty in threshold exceedance, for example through bootstrap p-values for each criterion, would make the checklist scores more principled.

Third, patient-facing uncertainty display is an open design problem. Conformal prediction sets at the individual level (CKD / not-CKD / ambiguous) produce clinically interpretable labels. How clinicians and patients respond to ambiguous predictions in a real consultation, and whether explicit uncertainty communication changes treatment decisions compared to point probability outputs, has not been tested in CKD care.
